\section{Global additive normalization and observable conditioning}
Proposition~\ref{prop:explicit} gives explicit degree and clock dependence
without claiming stable monomial interpolation in growing degree.
\begin{theorem}[Normalization, clock design, and spectral recovery]
\label{thm:main}
Suppose the true channels have full column rank. If $\ell\ge2d+1$
and $g=\gcd(f_1,\ldots,f_k)$ is monic, then
\begin{equation}\label{eq:kernel}
 \ker\cN_P=\{(q/g)s:\deg s<\deg g\},\qquad
 \dim\ker\cN_P=\deg g.
\end{equation}
For any design with $\ell\ge d+1$ and $\tau(P)>0$, every competitor
in $\cK_d$ satisfies
\begin{align}
 \|a-a'\|_2&\le C\min\{1,\delta/\tau(P)\},\label{eq:q-bound}\\
 d_\infty(\Lambda,\Lambda')&\le
 C\min\{1,[\delta(\kappa^{-1}+\tau(P)^{-1})]^{1/d}\}.
 \label{eq:poly-bound}
\end{align}
For simple roots within each true component, separated by $\nu>0$,
the spectral exponent is one, with the constant also depending on
$\nu$. Cross-component collisions are allowed. No rank or separation
hypothesis is imposed on a competitor.

At an interior parameter with internally simple roots, $\tau(P)=0$
gives a nonconstant local exact fibre with fixed channels and changing
aggregate roots. For every fixed admissible $k$, relatively prime
interior examples fail identification at $2d$ clocks. This worst-case
statement is realized by an affine-line matrix curve. If the component
polynomials span all polynomials of degree at most $d$, then $d+2$
clocks suffice and $d+1$ do not.
\end{theorem}
The proof, including the clock-complexity alternatives, is retained
in the appendices. To state its observable strengthening, reduce
$K$ in orthonormal frames of its rank-$k$ row and column spaces,
write $M(z)=\sum_{j=0}^dM_jz^j$, and put $A_0=M_d$.
For each distinct complete component polynomial $f$, let $E_f$ be
the joint spectral projector of the commuting $M_jA_0^{-1}$, and set
\[
 B_d(P)=\sum_f\|A_0^{-1}E_f\|_\op,
 \qquad \eta_d(P)=(B_d(P)+\tau(P)^{-1})^{-1}.
\]
Theorem~\ref{thm:joint} proves observability, frame independence, and
$d_\infty(\Lambda,\Lambda')\le
C\min\{1,[\delta/\eta_d(P)]^{1/d}\}$.
Separated clusters of at most $m$ roots per component have exponent
$1/m$. The two losses add; they are not multiplied.

In degree one, $f_b=z-\xi_b$, put
\begin{equation}\label{eq:affine}
 h=\sum_b\alpha_b\xi_b,\quad
 A=U\diag(\alpha)V^{\mathsf T},\quad
 C=U\diag(\alpha_b\xi_b)V^{\mathsf T},\quad B=hA-C.
\end{equation}
Then $P(T)=A+B/(T-h)$, $K(z)=zA-C$, and at three clocks
$\tau(P)\asymp\gamma(P)=\|B\|_F$.
The affine condition is $\sig=\kappa\gamma/(\kappa+\gamma)$,
zero on either boundary, with exact-fibre envelope
\begin{equation}\label{eq:envelope}
 \sig_*(P)=\max\{\sig(\omega):\mathcal P(\omega)=P\}.
\end{equation}
The residue condition is developed in Appendix~\ref{sec:residue};
$B_1$ is its residue sum and $\eta_1\asymp\eta$.
Appendix~\ref{sec:algorithm} gives the least-squares and polynomial-root
estimator without a supplied signal floor. This is exact-arithmetic
stability, not a bit-complexity or confidence-optimization assertion.
